\chapter{Harmonic analysis and spectral time series}
A continuous signal is a function $x:\R\to\R$ or $\C$.
The Fourier transform, when meaningful in the appropriate sense, is
$X(f)=\int_{\R}x(t)e^{-2\pi ift}\,dt$.
A length-$N$ recording is a different object: a discrete sequence.
An FFT of the recorded sequence does not directly reveal a unique
continuous generating signal without sampling assumptions.

\entry{57}{Euler's formula and Fourier orthogonality}
{$e^{i\theta}=\cos\theta+i\sin\theta$ and
$\int_0^T e^{2\pi i(k-\ell)t/T}\,dt=0$ for distinct integers
$k,\ell$.}
{Derive discrete spectral coefficients and phase representations.}
\entry{58}{Parseval's identity and Plancherel's theorem}
{For appropriate Fourier expansions/transforms,
energy in time equals energy in frequency, e.g.\
$\|x\|_2^2=\|X\|_2^2$ under a unitary convention.}
{Normalize power and check spectral energy.}
\entry{59}{Fourier convolution theorem}
{For suitable integrable functions, $\Fourier(f*g)=
\Fourier(f)\Fourier(g)$, with conventions fixed.}
{Relate a digital filter's impulse response to its transfer function.}
\entry{60}{Fourier shift and modulation identities}
{Translation of a signal multiplies its spectrum by a complex
phase; multiplication by a complex sinusoid translates frequency.}
{Understand lags, phase offsets, and analytic demodulation.}
\entry{61}{Heisenberg--Gabor uncertainty inequality}
{With standard centering and Fourier conventions, a nonzero
finite-energy waveform has time spread and frequency spread
whose product obeys a nonzero lower bound.}
{Understand the temporal and spectral resolution trade-off.}
\entry{62}{Paley--Wiener theorems}
{Appropriate compact-support and exponential-type conditions
characterize special classes of Fourier transforms;
exact simultaneous localization in time and frequency is impossible
for a nonzero ordinary waveform.}
{Explain support limitations of filters and short-time windows.}
\entry{63}{Riemann--Lebesgue lemma}
{The Fourier transform of an $L^1$ function vanishes
at infinity.}
{Understand high-frequency decay under integrability assumptions.}
\entry{64}{Bochner and Herglotz spectral representation}
{Continuous positive-definite functions on $\R$ are Fourier
transforms of finite nonnegative measures; discrete-time analogues
use nonnegative measures on the frequency circle.}
{Characterize admissible stationary covariance functions and spectra.}
\entry{65}{Cramer's spectral representation}
{A zero-mean weakly stationary second-order process can be
represented via orthogonal-increment random spectral measures
under the standard hypotheses.}
{Interpret stochastic LFP spectra rather than deterministic FFTs alone.}
\entry{66}{Wold decomposition}
{A purely nondeterministic covariance-stationary discrete-time
process has a one-sided linear innovations representation;
general stationary processes decompose into deterministic and
purely nondeterministic parts.}
{Understand assumptions behind linear innovations and VAR models.}
\entry{67}{Yule--Walker equations}
{For a stationary autoregressive process, lagged covariance
satisfies linear equations determined by AR coefficients.}
{Relate autocovariance to AR spectral estimators.}
\entry{68}{Wiener--Hopf normal equations}
{The least-squares optimal linear predictor has a residual
orthogonal to the predictor space, yielding covariance normal equations.}
{Derive linear decoding and linear prediction filters.}
\entry{69}{Spectral coherence bound}
{A positive-semidefinite two-signal spectral matrix satisfies
$|S_{xy}(f)|^2\leq S_{xx}(f)S_{yy}(f)$ where the spectra exist;
hence $0\leq C_{xy}(f)\leq1$ for positive denominators.}
{Interpret magnitude-squared LFP coherence.}
\entry{70}{Hilbert-transform Fourier multiplier identity}
{On $L^2(\R)$, with appropriate Fourier convention,
$\Fourier(\Hilbert x)(f)=-i\operatorname{sgn}(f)\Fourier x(f)$.}
{Build analytic signals and interpret instantaneous phase.}

\begin{definition}
For zero-mean jointly wide-sense stationary signals $x,y$,
let $R_{xy}(\tau)=\E[x(t+\tau)\overline{y(t)}]$.
Under appropriate integrability, $S_{xy}=\Fourier R_{xy}$.
Magnitude-squared coherence is
$C_{xy}=|S_{xy}|^2/(S_{xx}S_{yy})$ wherever the
denominator is strictly positive.
\end{definition}
\begin{caution}
Coherence is spectral dependence, not anatomical connectivity.
Common reference, volume conduction, leakage, sample-size bias,
and nonstationarity can distort interpretation
\cite{bastos2016,cohen2018}.
\end{caution}
